\section{Experimental design}\label{sec:design}

The study follows the ADEMP structure of Morris et al.~\cite{morris}: aims, data-generating mechanisms, estimands, methods
and performance measures, planned and stated before any run. The hypotheses were written in the study plan of the product
before the harness existed; they are restated here unchanged, and Section~\ref{sec:results} says for each whether it held.

\subsection{Aims and hypotheses}

The aim is to measure, on plants whose truth is known, when the calibrated predictor beats the empirical ones, and to check
each theoretical statement of Section~\ref{sec:theory} against simulation. Six hypotheses:
\begin{itemize}
\item[H1] (risk decomposition) The simulated mean squared error of the predictor with given scalars equals the risk of
Theorem~\ref{thm:risk} in every cell, within Monte Carlo error; at small shrinkage it beats least squares on $\mathbf Z$ if and
only if the curvature term of Corollary~\ref{cor:align} is positive, which on the plants with a scaled engine curvature means
a share $\alpha$ of that curvature above one half.
\item[H2] (capacity) On throughput with the engine's affine capacity, the gain over least squares on $\mathbf Z$ is zero up to
the shrinkage effect, for every design; with a capacity law curved like the plant's it grows with the width of the operating
window and is largest under extrapolation.
\item[H3] (setting-only line) Against the line in the setting, the crossing size $n^\star$ of Proposition~\ref{prop:line}
locates the number of surveys at which the method overtakes the line.
\item[H4] (priors) The three discrepancy priors change the predictive risk less than they change the meaning of the scalars:
their mean squared errors differ from one another by less than they differ from least squares.
\item[H5] (design) The throughput of surveys at a fixed speed cannot narrow the capacity base beyond the ceiling of
Proposition~\ref{prop:bound}; an added product-size output informs the population-balance scalars, and surveys at varying
speed inform the capacity base, by more than doubling the number of surveys at a fixed speed does.
\item[H6] (sampler) The sampler is calibrated: the ranks of simulation-based calibration are uniform for every parameter and for
the likelihood test quantity; and on the verification cells the full posterior's risk lies between those of the oracle and of
the linearized learned predictor, or within Monte Carlo error of the latter.
\end{itemize}

\subsection{Data-generating mechanisms}

Four virtual plants, each a mean function for throughput and power plus Gaussian noise with the standard deviations
$\sigma_{\mathrm{tph}}=\sigTph$~t/h and $\sigma_{\mathrm{kW}}=\sigKw$~kW, the posterior means of the fit to the ten
surveys~\cite{cb002}. The scalars $\boldsymbol\theta_0$ and the discrepancy coefficients $\mathbf d_0$ of the plants are the
posterior means of that fit (Table~\ref{tab:scalars}), and $\tilde{\mathbf z}_{\mathrm{ref}}$ are the regressors standardized
with the ten surveys' mean and deviation, fixed for the plant whatever the training design.
\begin{itemize}
\item V0, the well-specified plant: $\mu_t(\mathbf x)=f_t(\boldsymbol\theta_0;\mathbf x)+\tilde{\mathbf z}_{\mathrm{ref}}(\mathbf
x)^{\top}\mathbf d_{0,t}$, the model of~\eqref{eq:model} at fixed coefficients.
\item V1($\alpha$, $\kappa_q$), the plant with a controlled share of the engine's curvature in power: with $\mathbf b_0$ the
least-squares plane of $f_{\mathrm{kW}}(\boldsymbol\theta_0;\cdot)$ over the training window and $e_0$ the engine's departure
from it,
\begin{equation}
\mu_{\mathrm{kW}}(\mathbf x)=\tilde{\mathbf z}_{\mathrm{ref}}^{\top}(\mathbf b_0+\mathbf d_{0,\mathrm{kW}})+\alpha\,e_0(\mathbf x)+\kappa_q\,\|e_0\|\,q(\mathbf x),
\label{eq:v1}
\end{equation}
with $q$ the square of the standardized setting orthogonalized against $(1,z_c,z_f,e_0)$ over the window and scaled to unit
root mean square. The study runs $\alpha\in\{-1,0,0.5,1,2\}$ with $\kappa_q=0$, and $(\alpha,\kappa_q)=(1,2)$. V1(1, 0) is V0.
By Corollary~\ref{cor:align} the curvature gain is $(2\alpha-1)\|e_0\|^2$, positive exactly when $\alpha>1/2$, and the added
curvature $q$ leaves it unchanged while lowering the alignment.
\item V2, the plant with the published capacity law: throughput is the capacity of Rocha et al.~\cite{rocha} with their
published exponents and constant, $0.84\,Q_T(\rho_a/1600)(16/W_i)^{\eta}(100/f_{80})^{\epsilon}(c/38)^{\nu}$ with
$\nu=1.27$, curved in the setting and in the feed size; power is the engine's specific energy at $\boldsymbol\theta_0$ times
that throughput, $\zeta_0W_0(\mathbf x)\mu_{\mathrm{tph}}(\mathbf x)+P_{n,0}$, so it curves where the throughput does.
\item C-cap, the capacity study: the throughput of V2 predicted with three engines whose capacity is $q\,g(\mathbf x)$ with one
learned scale $q$: the shipped affine form $g=(c+17)/38$; the form of Rocha et al. itself; and the same form with its setting
and feed exponents halved, a curvature of the right sign and half the size.
\end{itemize}
For V2 and C-cap the oracle scalars are the $L_2$ projection of Tuo and Wu~\cite{tuowu}: the scalars that minimize the squared
distance between the engine and the plant over the training window, weighted by $1/\sigma_t^2$ across the outputs.

\subsection{Operating designs}

Training points are drawn from one of three windows: the surveys window, $c\in[38,55]$~mm and $f_{80}\in[47,140]$~mm, the
engine's calibrated envelope and the range of the ten surveys; the wide window, $c\in[25,70]$~mm and $f_{80}\in[40,140]$~mm with
$f_{80}\ge c+10$~mm; and the extrapolation design, trained on the surveys window and tested on the part of the wide window
outside it. The real design trains on the ten surveys themselves. The number of training surveys is $n\in\{10,20,40,80,160\}$.
Each window has a pool of 4,000 operating points on which the engine, its Jacobian and the plant are evaluated once; a
replicate draws its training points and 200 test points from the pools, and its noise. When training and test points come from
the same pool, the test points are drawn from the points not used for training. A design whose matrix $\mathbf G$ has a
condition number above $10^{10}$ is redrawn and counted; no design of the study needed it, and no pool point lost its engine
evaluation. Table~\ref{tab:cells} lists the cells.

\begin{table}[!t]
\caption{The cells of the study: each virtual plant on each design, with the number of training surveys, the predictors run
and the number of replicates; every cell uses 200 test points per replicate.}
\label{tab:cells}
\centering\footnotesize
\resizebox{\columnwidth}{!}{\begin{tabular}{@{}llllrr@{}}
\toprule
Plant & Outputs & Designs & n & Predictors & Cells \\
\midrule
V0 & throughput, power & real; surveys, wide, extrapolation & 10; 10 to 160 & 11 & 16 \\
V1($\alpha$, 0), 5 values & power & the same & the same & 11 & 80 \\
V1(1, 2) & power & the same & the same & 11 & 16 \\
V2 & throughput, power & the same & the same & 11 & 16 \\
C-cap, 3 engines & throughput & the same & the same & 9 & 48 \\
\bottomrule
\end{tabular}
}
\end{table}

\subsection{Estimands, methods and performance measures}

The estimand of a cell is the prediction risk~\eqref{eq:risk} of each predictor averaged over the test distribution of the
cell, and the coverage and the width of its 80\pct{} band. Eleven predictors run on every replicate: the method with the
oracle scalars and with the nominal scalars~\eqref{eq:cbmean}; the method with the scalars learned in the linearized model
under the three priors~\eqref{eq:linpred}; least squares on $\mathbf Z$; the setting-only line; the engine alone at the nominal
and at the oracle scalars; the engine refit without a discrepancy (the linearized maximum a posteriori of the scalars); and the
training mean. The regressions carry Student-$t$ prediction intervals with the residual variance; the Gaussian predictors carry
$\pm1.2816$ predictive standard deviations, with the plant's noise scale for the closed forms. Every performance number is a
mean over replicates with its Monte Carlo standard error, the standard deviation over replicates divided by the square root of
their number~\cite{morris}; the comparison of two predictors is the mean of the paired difference of their squared errors on
the same replicates, with the standard error of that difference. The coverage is the fraction of test points covered,
averaged over replicates. For the predictors with given scalars, each replicate also yields the risk of
Theorem~\ref{thm:risk} for its design and test points, and the cell reports the $z$ score of the simulated mean against the
theoretical one; the coverage of their bands is set beside the exact coverage of Corollary~\ref{cor:coverage}, which separates
the effect of the bias from that of the width as the bias-eliminated coverage of Morris et al.~\cite{morris} does.

\subsection{Replications and seeds}

Each cell has 1,900 replicates, the number that gives a Monte Carlo standard error of 0.5 points to a coverage near 95\pct{}
by Equation~(1) of Morris et al.~\cite{morris}, and below 1 point for an 80\pct{} coverage. The seed of a replicate is derived
from a fixed root, the cell's index and the replicate's index, so a result is reproducible and independent of the number of
processes; two runs of the same cell give identical files. The same seed fixes the design, the test points and the noise of a
replicate for every predictor, so the predictors are compared on identical data and the paired differences of
Section~\ref{sec:results} remove the variation that the design and the noise share. With 1,900 replicates of 200 test points the
Monte Carlo standard error of a mean squared error has a median of \mcseMedian\pct{} of its value over all cells and predictors,
at most \mcseMainMax\pct{} for the calibrated predictor, least squares and the line, and \mcseMax\pct{} for the heavy-tailed errors
of the orthogonal prior at ten surveys; that of a paired difference
is set by the spread of the difference, which is small where two predictors differ only by a shrinkage or a curvature term. The
full-sampler runs reuse the replicate seeds of the closed-form study, so their data are the same replicates with the same
indices.

\subsection{The full sampler}

Simulation-based calibration~\cite{talts} draws \sbcN{} parameter sets from the prior of the method (the eight scalars, the two
noise scales and the six discrepancy coefficients), simulates the ten surveys' design from each, fits the full posterior at
the reduced size of Section~\ref{sec:method} and records the rank of each true value among 99 evenly spaced posterior draws,
for $L+1=100$ possible ranks in 10 bins, which gives $N/B\ge20$ as Talts et al. recommend. The ranks of the marginal log
likelihood at the true parameters among its values at the draws are the data-dependent test quantity of Modr\'ak et
al.~\cite{modrak}, which detects a posterior equal to the prior. Uniformity is tested by a chi-square statistic per quantity.
The verification runs the full posterior on 50 replicates of three cells (V0 and V2 on the real design and V0 on the surveys
window with $n=10$), on exactly the data of the closed-form replicates with the same indices, and predicts 50 test points per
replicate from 200 posterior draws, so the full method and the closed forms are paired replicate by replicate.

\subsection{The design study}

The design study evaluates, in the linearized model at $\boldsymbol\theta_0$ with both outputs and their discrepancies, the
expected information gain~\eqref{eq:eig} and the ratio of posterior to prior standard deviation of each scalar for six designs:
D0 the ten surveys with throughput and power; D1 the same with the product size $P_{80}$ as a third output, with its own plane
discrepancy of 12\pct{} of its mean and a noise of 2, 5 or 10\pct{} of its mean; D2 the ten surveys with the speed set in turn
to 324, 360 and 396~rpm; D3 twenty surveys in the surveys window; D4 ten surveys in the wide window; and D5 the union of D1 and
D2. Random designs are summarized over 200 draws. The product size enters as a third output of the engine, the 80\pct{}
passing size of the predicted product, with its own plane discrepancy in the setting and the feed size; the three noise
levels bracket a measurement error that is not known for these surveys, whose product size distributions are printed only as
figures~\cite{rocha}. The engine and its
Jacobian at the speeds of D2 are evaluated at the same operating points as D0, so the two designs differ only in the speed.
Alongside, the identifiability analysis evaluates the ratio of the log capacity base for 10 to 640 surveys drawn from the surveys
window, at the fixed speed and with the speed varied, over 50 designs per size, and the floor~\eqref{eq:floor}.
